%-----------------------------------------------------------------------
\section{Connection to Generative Contact Mechanics}
\label{sec:connection}
%-----------------------------------------------------------------------

This section makes precise the relationship between the
singularity-theoretic framework developed in this volume and the
contact-geometric framework of \emph{Generative Contact Mechanics}
(Paper~1) and \emph{The dm3 Operator: Explicit Toy Model}
(Paper~2). We show that the operator sequence
$C \to K \to F \to U$
is the geometric precursor of the dm3 operator algebra, and that the
curvature threshold $\kappa^*$ corresponds canonically to the
embodiment threshold $\tau$.

%-----------------------------------------------------------------------
\subsection{From Trajectories to Limit Cycles}

In this volume the central object is a trajectory
$\gamma : [0,T] \to (X,g)$ undergoing a localized generative
transition driven by curvature, injectivity loss, and stabilization.

In GCM the central object is a hyperbolic limit cycle
$\Gamma \subset S$ embedded in a dissipative dynamical system
satisfying Axioms~1--8 of the dm3 framework.

The conceptual shift is:
\begin{itemize}[leftmargin=*]
  \item \emph{Principia Orthogona}: analyzes single transitions
        along trajectories.
  \item \emph{GCM}: studies persistent post-transition attractors
        and their algebraic interactions.
\end{itemize}

The mathematical link is that every admissible generative transition
induces a locally attracting invariant set, and conversely, every
dm3 limit cycle admits a neighborhood whose formation is governed by
a fold-type transition.

%-----------------------------------------------------------------------
\subsection{Curvature Threshold vs.\ Embodiment Threshold}

In this volume, folding is triggered when curvature reaches the
intrinsic geometric threshold
\[
  \kappa^*(x) = \frac{1}{\operatorname{foc}(x)},
\]
corrected by sectional curvature via the Rauch comparison theorem
(Section~\ref{sec:metric}).

In GCM, stochastic stability is governed by the embodiment threshold
\[
  \tau = \sqrt{c/\kappa},
\]
defined by the stochastic Lyapunov condition
$\mathcal{L}V \le -cV + \kappa\|\sigma\|^2$.

\begin{proposition}[Threshold Correspondence]\label{prop:threshold}
Let a generative transition produce a post-fold stabilized orbit
$\Gamma$ with Lyapunov function $V = d_g(\cdot,\Gamma)^2$. Then
the geometric condition $\kappa \uparrow \kappa^*$ is equivalent to
the dynamical condition that transverse contraction becomes dominant
over noise, yielding a finite embodiment threshold $\tau$.
\end{proposition}

\begin{proof}[Proof sketch]
As $\kappa \to \kappa^*$, the folding operator $F$ introduces
rank-1 loss and the post-fold unfolding $U$ selects a stable branch.
The resulting orbit $\Gamma$ has Floquet exponent
$\mu_{\max} < 0$ (transverse contraction). By the stochastic
Lyapunov condition of GCM, $\tau = \sqrt{c/\kappa}$ is finite
whenever $\mu_{\max} < 0$, i.e., precisely when curvature has
reached $\kappa^*$ and the fold has been completed. Thus
$\kappa^*$ is the geometric precursor of $\tau$: curvature
accumulation creates the conditions under which stochastic stability
becomes meaningful.
\end{proof}

%-----------------------------------------------------------------------
\subsection{Fold Map vs.\ dm3 Operator}

The folding operator $F$ in this volume is a rank-1 singular map:
$\operatorname{rank}(dF) = \dim(X_C) - 1$, producing a finite
branch set classified by $A_1$--$A_3$ singularities.

In GCM, the dm3 operator
$\mathcal{A}_{\mathrm{dm3}} = \phi^{T^*/4}$
is the time-$(T^*/4)$ map of a smooth flow on the contact manifold
$M = S \times \mathbb{R}$, acting discretely along a limit cycle.

\begin{proposition}[Fold--dm3 Correspondence]\label{prop:fold}
The fold map $F$ is the piecewise-smooth, pre-contact limit of the
dm3 operator in the following sense.
\begin{enumerate}[label=\textup{(\roman*)}]
  \item $F$ enforces non-injectivity geometrically via rank loss;
        $\mathcal{A}_{\mathrm{dm3}}$ enforces recurrence dynamically
        via orbital stability.
  \item Under the contact extension $M = X \times \mathbb{R}$
        \textup{(}the construction of GCM, \S6\textup{)}, the
        impulsive momentum jump at the fold
        $p(s_0^+) - p(s_0^-) = \mu\,n(s_0)$
        corresponds to the contact Hamiltonian correction
        $H_{\mathrm{diss}} = -\gamma V e^{-\beta z}$,
        which drives trajectories toward $\Gamma$ off the orbit.
\end{enumerate}
\end{proposition}

\begin{proof}[Proof sketch]
Both the momentum jump and $H_{\mathrm{diss}}$ are corrections
that vanish on the invariant set ($\Gamma$ or the post-fold branch)
and are nonzero off it. The momentum jump is a distributional
first-order correction in the Hamiltonian sense; $H_{\mathrm{diss}}$
is its smooth contact-geometric regularization as $\mu \to 0$ and
the contact manifold structure is imposed.
\end{proof}

%-----------------------------------------------------------------------
\subsection{Operator Grammar Correspondence}

The operator sequence of this volume maps onto the reduced dm3
operator grammar as follows.

\begin{center}
\begin{tabular}{ll}
\toprule
\emph{Principia Orthogona} & \emph{GCM / dm3} \\
\midrule
Compression $C$ & Basin contraction via $g_1$ \\
Curvature flow $K$ & Lyapunov descent $\dot{V} \le -cV$ \\
Fold $F$ & Contact bifurcation (Hopf, SN, NS) \\
Unfolding $U$ & Gradient flow to $\Gamma_{\mathrm{syn}}$ via $U_3$ \\
\bottomrule
\end{tabular}
\end{center}

The full dm3 grammar $g \to L \to R \to U$ extends this sequence by
adding multi-orbit coherence, resonance detection, and categorical
unification, which are intentionally absent from the present volume.

%-----------------------------------------------------------------------
\subsection{Singularity Classification vs.\ Bifurcation Diagram}

This volume proves that admissible generative transitions are
precisely the Whitney singularities $A_1$, $A_2$, $A_3$.
Paper~2 proves that dm3 systems undergo exactly four bifurcations:
contact Hopf, saddle-node of limit cycles, Neimark--Sacker, and
slow-fast crossover.

\begin{remark}[Classification Correspondence]
The $A_1$--$A_3$ singularities classify the local geometry of dm3
bifurcations under the projection $M = S \times \mathbb{R} \to S$.
The fold $A_1$ corresponds to the contact Hopf and saddle-node
bifurcations (local rank loss in the radial direction); the cusp
$A_2$ corresponds to the Neimark--Sacker (loss of a second
direction at resonance); the swallowtail $A_3$ to the slow-fast
crossover (third-order degeneracy in the $z$-direction). Higher
singularities are excluded in both frameworks: here by the Morse
condition on $\Phi$; in dm3 by contact normal form rigidity
(Theorem~C of Paper~1).
\end{remark}

%-----------------------------------------------------------------------
\subsection{Summary and Interpretation}

The three papers in this series have the following relationship.

\begin{itemize}[leftmargin=*]
  \item \textbf{This volume}: explains why folds must occur and
        classifies their local geometry.
  \item \textbf{Paper~1 (GCM)}: explains what stable structures
        exist after they occur, in the contact-geometric setting.
  \item \textbf{Paper~2 (dm3)}: proves that the resulting dynamics
        are computable, stable, and structurally closed, by explicit
        global analysis of the toy model.
\end{itemize}

No claim is made that every dm3 system arises from a literal
curvature fold in a physical manifold. The precise claim is:

\begin{quote}
Every dm3 system admits a local geometric model whose generative
event is $\mathcal{A}$-equivalent to a fold-type transition governed
by a curvature threshold $\kappa^*$.
\end{quote}

This claim is supported by
Propositions~\ref{prop:threshold} and~\ref{prop:fold}, which
establish the threshold correspondence and the fold--dm3
correspondence respectively.
